\documentclass[11pt]{article}
\usepackage[a4paper,margin=1in]{geometry}
\usepackage{amsmath,amssymb,amsthm,booktabs,hyperref}
\newtheorem{theorem}{Theorem}
\newtheorem{lemma}{Lemma}
\newcommand{\ind}{\mathbf{1}}
\title{An odd-divisor inverse for a shifted Lambert factorization\\
\large A counterexample and corrected formula for MathDB 356220}
\author{Principia Math In Progress}
\date{18 September 2026}
\begin{document}
\maketitle
\begin{abstract}
The proposed inverse in Schmidt's Conjecture 2.26 fails at $(n,k)=(58,1)$:
it gives $613766$, whereas the inverse entry is $613767$.
Ordinary M\"obius inversion on odd integers gives a corrected formula for
every entry and every invertible formal prefactor.
This is an authoring-agent checked manuscript, without independent review.
Campaign credit is zero pending inspection of the related final journal article.
\end{abstract}

\section*{Task 1. Refutation and complete corrected inverse}
Let $p(j)$ be the ordinary partition number, with $p(0)=1$ and
$p(j)=0$ for negative integers $j$. Define $S=(s_{n,k})_{n,k\geq1}$ by
\begin{equation}\label{factor}
 \sum_{k\geq1}\frac{f(k)q^k}{1-q^{2k+1}}
 =\frac1{(q;q)_\infty}
   \sum_{n\geq1}\sum_{k=1}^n s_{n,k}f(k)q^n,
 \qquad (q;q)_\infty=\prod_{j\geq1}(1-q^j).
\end{equation}
These are formal power series; every coefficient uses finitely many terms.
The matrix $S$ is unit lower triangular, and $s^{(-1)}_{n,k}$ denotes
the entry of its inverse.

Conjecture 2.26 of the thesis \cite{thesis}, printed page 31, proposes
\begin{align}
 s^{(-1)}_{n,k}\stackrel{?}{={}}&
 p(n-k)-\sum_{\substack{1\leq i\leq n\\n\equiv i\ (2i+1)}}
 p\left(\frac{n-i}{2i+1}-k\right)\label{proposed}\\
 &+\sum_{m=2}^n
 \sum_{\substack{1\leq i\leq n\\
 n\equiv p(m+1)i+p(m-1)\ (p(m+1)(2i+1))}}
 p\left(\frac{n-p(m+1)i-p(m-1)}{p(m+1)(2i+1)}-k\right).\notag
\end{align}
The congruences in the summation conditions have the same meaning as the
source's Iverson factors. This formula also appears as equation (5) in
Remark 2.7 of \cite{preprint}. The function $p$ in the source is the
partition function throughout.

\begin{theorem}\label{main}
For all integers $n\geq k\geq1$, the actual inverse is
\begin{equation}\label{inverse}
 \boxed{\displaystyle
 s^{(-1)}_{n,k}
 =\sum_{d\mid 2n+1}\mu(d)\,
 p\left(\frac{(2n+1)/d-1}{2}-k\right).}
\end{equation}
Here $\mu$ is the classical M\"obius function.
In particular $s^{(-1)}_{58,1}=613767$, while the right side of
\eqref{proposed} is $613766$. Thus the conjecture is false.
\end{theorem}

\subsection*{Inversion of the odd-divisor kernel}
Write $C(q)=\sum_{j\geq0}c_jq^j$, where $c_0=1$, and
$C(q)^{-1}=\sum_{j\geq0}r_jq^j$, with $r_j=0$ when $j<0$.
Replacing $(q;q)_\infty$ by $C(q)$ in \eqref{factor} defines a matrix $S_C$.
Let $T_C=(c_{n-k})_{n,k\geq1}$, so that
$T_C^{-1}=(r_{n-k})_{n,k\geq1}$.

\begin{lemma}\label{odd}
The coefficient matrix of $q^k/(1-q^{2k+1})$ is
\[
 D_{n,k}=\ind_{\,2k+1\mid 2n+1}\qquad(n,k\geq1).
\]
Its inverse is
\[
 E_{n,k}=
 \begin{cases}
 \mu\bigl((2n+1)/(2k+1)\bigr),&2k+1\mid2n+1,\\
 0,&\text{otherwise}.
 \end{cases}
\]
\end{lemma}
\begin{proof}
The coefficient in question is $1$ precisely when
$n=k+\ell(2k+1)$ for an integer $\ell\geq0$.
This is equivalent to $2n+1=(2\ell+1)(2k+1)$.
Conversely, the quotient of two odd integers is odd, so divisibility gives
such an $\ell$.

If $2k+1$ does not divide $2n+1$, the $(n,k)$ entries of both $DE$ and
$ED$ vanish. Otherwise put $M=(2n+1)/(2k+1)$. The intermediate odd
integers $2j+1$ are exactly $(2k+1)e$ with $e\mid M$.
Consequently each product entry is $\sum_{e\mid M}\mu(e)$, which is
$1$ for $n=k$ and $0$ otherwise. The divisor identity follows directly
from $\prod_{\ell\mid M,\ \ell\ {\rm prime}}(1-1)$ when $M>1$.
There is no missing index at the odd integer $1$: every intermediate
integer is at least $2k+1\geq3$.
\end{proof}

\begin{proof}[Proof of the all-entry formula]
Multiplication by $C(q)$ gives $S_C=T_CD$. Thus
\[
 S_C^{-1}=ET_C^{-1},\qquad
 (S_C^{-1})_{n,k}
 =\sum_{\substack{k\leq j\leq n\\2j+1\mid2n+1}}
 \mu\left(\frac{2n+1}{2j+1}\right)r_{j-k}.
\]
Substitute $d=(2n+1)/(2j+1)$. This proves
\begin{equation}\label{general}
 (S_C^{-1})_{n,k}
 =\sum_{d\mid2n+1}\mu(d)\,
 r_{\frac{(2n+1)/d-1}{2}-k}.
\end{equation}
All omitted indices have negative subscripts and therefore contribute zero.
For $C(q)=(q;q)_\infty$ we have $r_j=p(j)$, proving \eqref{inverse}.
All matrix products are finite at each entry. In particular, inversion of
the leading $N$ by $N$ block agrees with this formula whenever $n\leq N$.
\end{proof}

\subsection*{Exact evaluation at $(58,1)$}
Since $117=3^2\cdot13$, the nonzero M\"obius terms have
$d=1,3,13,39$. Equation \eqref{inverse} gives
\[
 p(57)-p(18)-p(3)+p(0)
 =614154-385-3+1=613767.
\]
For the negative sum in \eqref{proposed}, divisibility gives the following
and only the following nonzero terms:
\[
\begin{array}{c|rrrr}
i&1&4&6&19\\ \hline
(58-i)/(2i+1)-1&18&5&3&0
\end{array}
\]
For the positive double sum, a nonzero term requires
\[
58\geq p(m+1)(3i+1)+p(m-1).
\]
For $m\geq6$ the right side is at least $4p(7)+p(5)=67$.
For $m=2,3,4,5$, the possible $i$ from this bound are respectively
$1,\ldots,6$; $1,\ldots,3$; $1,2$; and $1$.
Testing the indicated congruences leaves just $(m,i)=(2,1),(2,6)$,
with partition indices $5,0$. Hence the proposed value is
\[
 p(57)-\bigl(p(18)+p(5)+p(3)+p(0)\bigr)
       +\bigl(p(5)+p(0)\bigr)
 =614154-396+8=613766.
\]
This is a failure of the mathematical formula repeated in both sources.
The corrected theorem answers the inverse-matrix question for all entries.

\subsection*{Verification and source limits}
The accompanying standard-library Python script builds $S$ directly from
finite products and geometric series, then inverts it by exact triangular
recursion. It checks all $1711$ entries through row $58$, including the
first failure of \eqref{proposed}; it also tests \eqref{general} for two
other prefactors. These finite tests supplement the proof above.
Classical M\"obius inversion and formal-series matrix multiplication are
credited standard tools, not new general theorems.

The complete 2022 thesis and 2017 arXiv preprint were inspected.
The final journal version of \cite{preprint}, published in 2020,
was not accessible in full. The literature and priority check therefore
remains pending, with zero campaign credit. This limitation concerns
eligibility, not the exact discrepancy for the displayed formula.

\section*{Four proposed follow-up tasks}
\paragraph{Task 2. General integer slopes.}
For each integer $\alpha\geq3$, obtain an explicit arithmetic inverse for
the coefficient matrix of $q^k/(1-q^{\alpha k+1})$.
Identify the appropriate divisor-poset M\"obius function and its support.
\paragraph{Task 3. The weighted denominator.}
Determine the inverse entries for $q^k/(1-zq^{2k+1})$ as polynomials in $z$,
and give a factorization-chain interpretation for their coefficients.
\paragraph{Task 4. Positivity of the corrected entries.}
Determine whether the numbers in \eqref{inverse} are nonnegative for every
$n\geq k\geq1$. If they are, seek a partition interpretation of the
odd-divisor inclusion-exclusion sum.
\paragraph{Task 5. Classify the erroneous formula's agreement set.}
Characterize the pairs $(n,k)$ for which \eqref{proposed} equals
\eqref{inverse}, and determine the density or growth of rows containing
at least one disagreement.

\begin{thebibliography}{9}
\bibitem{thesis} M. D. Schmidt,
\emph{Factorization theorems and canonical representations for generating
functions of special sums}, Ph.D. thesis, Georgia Institute of Technology,
2022, Conjecture 2.26, p.~31.
\href{https://arxiv.org/abs/2209.12287}{arXiv:2209.12287}.
\bibitem{preprint} M. Merca and M. D. Schmidt,
\emph{Factorization theorems for generalized Lambert series and applications},
\href{https://arxiv.org/abs/1712.00611}{arXiv:1712.00611}, Remark 2.7,
equation (5), pp.~8--9.
Published as \emph{Ramanujan J.} \textbf{51} (2020), 391--419,
\href{https://doi.org/10.1007/s11139-018-0095-7}{doi:10.1007/s11139-018-0095-7}.
The full published version was not inspected.
\end{thebibliography}
\end{document}
